\documentclass[11pt]{article}
\usepackage[margin=1in]{geometry}
\usepackage{amsmath,amssymb,amsthm}
\usepackage[T1]{fontenc}
\usepackage{lmodern}
\usepackage{listings}
\lstset{basicstyle=\ttfamily\small,columns=fullflexible,keepspaces=true,
  literate={Ω}{{$\Omega$}}1 {∃}{{$\exists$}}1 {→}{{$\to$}}1
  {↔}{{$\leftrightarrow$}}1 {ℕ}{{$\mathbb N$}}1 {∧}{{$\land$}}1}
\usepackage{microtype}
\usepackage{hyperref}
\hypersetup{colorlinks=true,linkcolor=blue,citecolor=blue,urlcolor=blue,
  pdftitle={A classification of solvable spinning-switches puzzles},pdfauthor={P. Graham}}
\newtheorem{theorem}{Theorem}
\newtheorem{lemma}[theorem]{Lemma}
\newcommand{\K}{G^{\Omega}}
\newcommand{\al}{\alpha}
\title{A classification of solvable spinning-switches puzzles}
\author{P. Graham\\ \small\url{https://orcid.org/0009-0007-8493-1505}}
\date{October 31, 2026}
\pagestyle{plain}
\begin{document}
\maketitle
\begin{abstract}
Spinning-switches puzzles ask for a predetermined sequence of moves that
works despite hidden permutations of the switches. Kagey formulated the
problem using a finite group of switch states and a permutation group of
positions, proved solvability when both groups are $p$-groups for one
prime $p$, and classified the case of abelian switch groups. We prove the
converse without an abelian hypothesis: for nontrivial finite groups $G$
and $H$, with $H$ acting faithfully on a finite set $\Omega$, a surjective
strategy for $G\wr_\Omega H$ exists exactly when $G$ and $H$ are $p$-groups
for the same prime. This answers Kagey's Open Question 31. The necessity
proof restricts the adversary to a prime-order element and tracks excluded
configurations by cosets of prime-power subgroups. We also give a central-series
proof of the known sufficient condition. The classification and its
consequences for $S_3\wr C_2$ and $A_4\wr C_3$ are formalized in Lean.

This work has not yet been reviewed by experts.
The proofs, the Lean formalization and the text were produced by AI models (from OpenAI and Anthropic), working under the author's direction.
\end{abstract}
\medskip
\noindent\textbf{MSC 2020:} 20D15, 05E18.\quad
\textbf{Keywords:} spinning switches, wreath products, finite groups.

\section{Introduction}
A player changes hidden switches while an adversary permutes their positions.
The player must choose the moves in advance. Kagey's model~\cite{Kagey}
assigns a group $G$ to each switch and lets a group $H$ act on the set of
positions $\Omega$. Write $K=\K$ for the configuration group, with
pointwise multiplication, and put
\[
  (\al_hx)(\omega)=x(h^{-1}\omega),\qquad
  W=K\rtimes_\al H,\qquad
  (x,h)(y,a)=(x\al_h(y),ha).
\]
A move is an arbitrary element of $K$; in particular, it may change several
switches at once. For fixed moves $k_1,\ldots,k_N$ and adversary choices
$h_1,\ldots,h_N$, put $P_0=1$ and
$P_j=(k_1,h_1)\cdots(k_j,h_j)$. Write $P_j=(b_j,a_j)$.
Following Definition 9 of~\cite[arXiv v1]{Kagey}, the moves form a
\emph{surjective strategy} if, for every adversary sequence,
$\{b_j:0\leq j\leq N\}=K$. Equivalently, from every initial configuration
$x_0$, some $x_0b_j$ equals $1$. No observation of the spins is available.

\begin{theorem}\label{thm:main}
Let $G$ and $H$ be nontrivial finite groups, and let $H$ act faithfully on
a finite set $\Omega$. There is a surjective strategy for $G\wr_\Omega H$
if and only if there is a prime $p$ for which both $G$ and $H$ are $p$-groups.
\end{theorem}

The action need not be transitive. A $p$-group here means a finite group
whose order is a power of $p$. The nontriviality assumptions matter:
trivial switches need no moves, and a trivial adversary permits a tour of
any finite configuration group. If $N$ is the kernel of the action, each
$b_j$ depends only on the images of the $h_i$ in $H/N$, so
$G\wr_\Omega H$ has a surjective strategy if and only if
$G\wr_\Omega(H/N)$ does. Hence in general a surjective strategy exists
if and only if $G=1$, or $H$ acts trivially, or $G$ and $H/N$ are
$p$-groups for one prime $p$.

All statement numbers below refer to arXiv:2210.09408v1, posted
October 17, 2022, of Kagey's 2023 paper~\cite{Kagey}.
Theorem 22 gives the sufficient condition in Theorem~\ref{thm:main}.
Theorem 23 gives the equivalence when $G$ is abelian.
Open Question 31 asks for the finite wreath products admitting a
surjective strategy. Our contribution is necessity for arbitrary finite
switch groups, including nonabelian groups; Section~\ref{sec:suff}
reproves the known sufficiency for a complete formal proof.
Earlier rotating-table and blind-bartender puzzles were studied by
Bar Yehuda, Etzion and Moran~\cite{BYEM} and Ehrenborg and
Skinner~\cite{ES}; Rabinovich's linear-action criterion~\cite{Rabinovich},
stated as Theorem 13 of~\cite[arXiv v1]{Kagey}, says that a faithful
linear action on a finite-field vector space of characteristic $p$
admits a strategy exactly when the acting group is a $p$-group.
That linear setting does not cover nonabelian switch groups; our coset
invariant supplies necessity in this case.

Theorem 26 of that version asserts solvability for $G\wr C_2$ when $G$
is generated by involutions. This assertion was already corrected:
ndkrempel reported on MathOverflow on December 7, 2025 that Theorem 26
fails, and Kagey acknowledged the error there on December 8, 2025,
noting that Example 27 fails~\cite{MO}.
The classification, not that correction, is the claim of this paper.
For example, it implies that $S_3\wr C_2$, with $C_2$ interchanging two
positions, is unsolvable, refuting Conjecture 32 (already implicit in
the Theorem 26 correction). It refutes Conjecture 33 for every $n\geq4$:
$A_n\wr C_3$ on three cyclically permuted positions is unsolvable
because $|A_n|$ is even; the $n=3$ case holds by Theorem 22.
For a nontrivial $p$-group $H$ acting faithfully, it settles Conjecture 34
exactly: a strategy exists precisely when $G$ is a $p$-group; the trivial
adversary case is always solvable, as above. The formalized negative
instances are $S_3\wr C_2$ and $A_4\wr C_3$, with group orders $6,2$
and $12,3$, respectively.

\section{The prime-order obstruction}
Fix a prime $q$, an element $h\in H$ of order $q$, and a position $\omega$
that $h$ moves. Its orbit $O$ has $q$ elements. Restrict the adversary to
$\langle h\rangle$. For $k\in K$ and $E\subseteq K$, define
\begin{equation}\label{eq:step}
 T_k(E)=\{1\}\cup
 \{y\in K:(a\mapsto y(h^na))k^{-1}\in E\text{ for every }n\geq0\}.
\end{equation}
For a move sequence, set $E_0=\{1\}$ and $E_i=T_{k_i}(E_{i-1})$.
These are configurations excluded from surviving play, expressed in the
adversary's moving coordinates.

\begin{lemma}\label{lem:back}
If $E_N\neq K$, the moves $k_1,\ldots,k_N$ are not a surjective strategy.
\end{lemma}
\begin{proof}
Choose $z_N\notin E_N$. Recursively,~\eqref{eq:step} gives $n_i$ such that
$z_{i-1}=(a\mapsto z_i(h^{n_i}a))k_i^{-1}\notin E_{i-1}$.
Every $z_i$ is nonidentity. Set $h_i=h^{-n_i}$ and $x_0=z_0$.
The multiplication in $W$ gives, inductively,
\begin{equation}\label{eq:moving}
\begin{aligned}
 x_0b_i &= x_0b_{i-1}\al_{a_{i-1}}(k_i)
          = \al_{a_{i-1}}(z_{i-1}k_i),\\
 \al_{a_i^{-1}}(x_0b_i)
        &= \al_{h_i^{-1}}(z_{i-1}k_i)=z_i,
\end{aligned}
\end{equation}
since $a_i=a_{i-1}h_i$. Thus $x_0b_i\neq1$ for every prefix,
so the base projections miss $x_0^{-1}$.
\end{proof}

Say that $y\in K$ is \emph{$P$-aligned on $O$}, for a subgroup $P\leq G$,
if $y(a)^{-1}y(b)\in P$ for all $a,b\in O$. This means its values on $O$
lie in one left coset of $P$.

\begin{lemma}\label{lem:coset}
If every member of $E$ is $P$-aligned on $O$ for a $q$-subgroup $P\leq G$,
then there is a $q$-subgroup $P'\leq G$ such that every member of
$T_k(E)$ is $P'$-aligned on $O$.
\end{lemma}
\begin{proof}
For $y\in T_k(E)$ with $y\neq1$, rotate any $t\in O$ to any $a\in O$
in~\eqref{eq:step} and apply alignment at $t,ht$. This gives
\begin{equation}\label{eq:diff}
 k(t)\bigl(y(a)^{-1}y(ha)\bigr)k(ht)^{-1}\in P
 \qquad(a,t\in O).
\end{equation}
Put
\[
 R=\bigcap_{t\in O}k(t)^{-1}Pk(t),\qquad
 C=\bigcap_{t\in O}k(t)^{-1}Pk(ht).
\]
If $T_k(E)=\{1\}$, take $P'=1$. Otherwise~\eqref{eq:diff} supplies
$c\in C$. We have $C=Rc$: indeed, $d\in C$ exactly when
$d c^{-1}\in R$. Also $c$ normalizes $R$. To see this, write
$v_t=k(t)c k(ht)^{-1}\in P$; for $r\in R$,
\[
 k(t)c r c^{-1}k(t)^{-1}
   =v_t\bigl(k(ht)r k(ht)^{-1}\bigr)v_t^{-1}\in P.
\]
This proves $cRc^{-1}\subseteq R$, hence equality by finiteness.
Multiplying $v_t v_{ht}\cdots v_{h^{q-1}t}$ shows
$k(t)c^q k(t)^{-1}\in P$, so $c^q\in R$.

The subgroup $R$ is a $q$-group, being contained in a conjugate of $P$.
Consequently $P'=\langle R,c\rangle=R\langle c\rangle$ is a $q$-group:
$R$ is normal in $P'$ and $P'/R$ is cyclic of order dividing $q$.
Equation~\eqref{eq:diff} puts every consecutive difference
$y(a)^{-1}y(ha)$ in $C\subseteq P'$. Multiplying consecutive differences
around $O$ gives $y(a)^{-1}y(b)\in P'$ for all $a,b\in O$.
The identity configuration satisfies the same conclusion.
\end{proof}

\begin{proof}[Necessity in Theorem~\ref{thm:main}]
Suppose a surjective strategy exists. For any prime $q$ dividing $|H|$,
Cauchy's theorem supplies an element $h$ of order $q$; faithfulness supplies
a position it moves. Starting with the trivial subgroup and $E_0=\{1\}$,
Lemma~\ref{lem:coset} gives a $q$-subgroup $P_N$ such that all of $E_N$
is $P_N$-aligned on $O$. Lemma~\ref{lem:back} forces $E_N=K$.
For each $g\in G$, choose a configuration with value $1$ at $\omega$ and
$g$ at $h\omega$. Alignment forces $g\in P_N$. Thus $G=P_N$ is a
$q$-group.

Choose a prime $p\mid |H|$, possible since $H\neq1$. Then $G$ is a
$p$-group. Because $G\neq1$, it cannot also be a $q$-group for a different
prime $q$. Every prime divisor of $|H|$ is therefore $p$, and $H$ is a
$p$-group as well.
\end{proof}

\section{Sufficiency and the formal theorem}\label{sec:suff}
Here is a self-contained proof of Kagey's sufficient condition. For the
full adversary, replace~\eqref{eq:step} by
\[
 T^H_k(E)=\{1\}\cup\{y:\al_a(y)k^{-1}\in E\text{ for all }a\in H\}.
\]
Let $E^H(s)$ be the result of applying these operations to $\{1\}$ along
a sequence $s$. If a play from $x_0$ never reaches $1$, its moving
configuration $z_j=\al_{a_j^{-1}}(x_0b_j)$ stays outside the corresponding
excluded set: using~\eqref{eq:moving} with $a=h_j$ in the universal
condition gives $\al_{h_j}(z_j)k_j^{-1}=z_{j-1}$, so induction excludes
$z_j$ from $T^H_{k_j}(E^H(k_1,\ldots,k_{j-1}))$.
Hence $E^H(s)=K$ implies that $s$ is surjective.

\begin{lemma}\label{lem:extension}
Let $M\lhd K$ be $H$-invariant. Suppose a sequence $s$ of moves in $M$
satisfies $M\subseteq E^H(s)$. If $u\in K\setminus M$ commutes modulo $M$
with every element of $W$, the same assertion holds for the strictly
larger normal, $H$-invariant subgroup $L=M\langle u\rangle$.
\end{lemma}
\begin{proof}
The commutator hypothesis makes $L$ normal in $K$ and $H$-invariant,
and makes $H$ fix each coset $Mu^j$. Form a sequence consisting of $r$
copies of $s$ separated by the move $u^{-1}$, where $r$ is the order of $u$.
After $i+1$ copies, the excluded set contains
$M\cup Mu^{-1}\cup\cdots\cup Mu^{-i}$. This follows by induction.
The separating move sends each guaranteed coset $Mu^{-j}$ to
$Mu^{-(j+1)}$ in the recurrence. Each subsequent move in $s$ preserves
these shifted cosets, because they are $H$-invariant and closed under
right multiplication by $M$. Meanwhile the block $s$ restores $M$:
every excluded set contains $1$, and the recurrence is monotone in $E$.
After $r$ copies all of $L$ is excluded. All moves used lie in $L$.
\end{proof}

\begin{proof}[Sufficiency in Theorem~\ref{thm:main}]
If $G$ and $H$ are $p$-groups, then $W$ is a finite $p$-group and so is
nilpotent. Begin with $M=1$ and the empty sequence. For $M\neq K$, let
$Z_i(W)$ be the upper central series, and choose the least $i$ with
$K\cap Z_i(W)\nsubseteq M$. Choose $u$ in this difference. Its commutators
with $W$ lie in $K\cap Z_{i-1}(W)\subseteq M$.
Lemma~\ref{lem:extension} enlarges $M$ while preserving the sequence
property. Finiteness makes this process reach $M=K$, giving a sequence
with $E^H(s)=K$, hence a surjective strategy.
\end{proof}

The accompanying Lean project proves the full statement, in namespace
\texttt{SpinningSwitches}, with the following parameters and type:
\begin{lstlisting}
variable {G H Ω : Type*} [Group G] [Group H] [MulAction H Ω]
theorem classification [Finite G] [Finite H] [Finite Ω]
    [FaithfulSMul H Ω] [Nontrivial G] [Nontrivial H] :
    (∃ ks : List (Ω → G), IsSurjectiveStrategy H ks) ↔
      ∃ p : ℕ, p.Prime ∧ IsPGroup p G ∧ IsPGroup p H
\end{lstlisting}
Here \texttt{IsSurjectiveStrategy} quantifies over every equally long list
of adversary moves and every base configuration, requiring a prefix
index from zero through the list length whose base projection equals
that configuration. It is exactly the definition in Section 1.
For finite groups, Mathlib's \texttt{IsPGroup} is equivalent to having
prime-power order. The development also checks the two unsolvable
examples stated above. It does not assert an optimal strategy length.

To reproduce the check, install elan, Git, and a C compiler, enter
\texttt{lean/} in the accompanying bundle, and run
\texttt{lake exe cache get}, then \texttt{lake build}.
The project pins Lean 4.34.0 and its Mathlib dependencies. The build checks
the proof, its statement type, small examples, and axiom guards.
The theorem uses only Lean's standard axioms \texttt{propext},
\texttt{Classical.choice}, and \texttt{Quot.sound}; no compiled evaluation
is used. This relies on Lean's kernel and the stated formal definitions.
The formalization uses Mathlib's finite-group and nilpotency
libraries~\cite{Mathlib}. A cached build takes several minutes.

\begin{thebibliography}{9}
\bibitem{Kagey}
P. Kagey, \emph{Spinning switches on a wreath product},
Journal of Combinatorial Theory, Series A \textbf{200} (2023), 105795.
\href{https://doi.org/10.1016/j.jcta.2023.105795}{doi:10.1016/j.jcta.2023.105795}.
Statement numbers refer to
\href{https://arxiv.org/abs/2210.09408v1}{arXiv:2210.09408v1}, October 17, 2022.
\bibitem{MO}
ndkrempel and P. Kagey, comments of December 7 and December 8, 2025,
on Kagey's answer to \emph{Game on groups (generalization of spinning switches puzzle)},
MathOverflow, \url{https://mathoverflow.net/a/423643}.
\bibitem{BYEM}
R. Bar Yehuda, T. Etzion and S. Moran,
\emph{Rotating-table games and derivatives of words},
Theoretical Computer Science \textbf{108} (1993), 311--329.
\href{https://doi.org/10.1016/0304-3975(93)90196-Z}{doi:10.1016/0304-3975(93)90196-Z}.
\bibitem{ES}
R. Ehrenborg and C. M. Skinner, \emph{The blind bartender's problem},
Journal of Combinatorial Theory, Series A \textbf{70} (1995), 249--266.
\href{https://doi.org/10.1016/0097-3165(95)90092-6}{doi:10.1016/0097-3165(95)90092-6}.
\bibitem{Rabinovich}
Y. Rabinovich, \emph{A generalization of the Blind Rotating Table game},
Information Processing Letters \textbf{176} (2022), 106233.
\href{https://doi.org/10.1016/j.ipl.2021.106233}{doi:10.1016/j.ipl.2021.106233}.
\bibitem{Mathlib}
The mathlib Community, \emph{The Lean mathematical library},
Proceedings of CPP 2020, 367--381.
\href{https://doi.org/10.1145/3372885.3373824}{doi:10.1145/3372885.3373824}.
\end{thebibliography}
\end{document}
